\section{\label{sec:defns}Light-front and quasidistributions}

Throughout this work, I focus on flavor nonsinglet unpolarized quasi and light-front PDFs. The extension to polarized quasi-PDFs is straightforward and much of the discussion applies equally to pseudodistributions. I do not consider the flavor singlet case, which introduces additional mixing with the gluon distribution; this mixing is studied for the case of the quasidistribution at one loop in perturbation theory in Ref.~\cite{Wang:2017qyg}.

\subsection{Light-front PDFs}

I denote renormalized light-front PDFs by $f(\xi,\mu)$, where the renormalization scale is $\mu$, and define $\xi=k^+/P^+$. Here the light-front coordinates are $(x^+,x^-,\boldsymbol{x}_\mathrm{T})$ such that $x^\pm = (t\pm z)/\sqrt{2}$. In general, one can write the renormalized light-front PDFs in terms of the bare light-front PDFs, $f^{(0)}(\xi)$, as
\begin{equation}
f(\xi,\mu) = \int_\xi^1 \frac{\mathrm{d}\zeta}{\zeta}{\cal
Z}\left(\frac{\xi}{\zeta},\mu\right)f^{(0)}(\zeta).
\end{equation}
The bare PDF is defined as \cite{Collins:2011zzd}
\begin{align}
 {} &f^{(0)}(\xi) = \int_{-\infty}^\infty
\frac{\mathrm{d}\omega^-}{4\pi}
e^{-i\xi P^+\omega^-} \nonumber\\
{} & \quad \times \left\langle P \left|
T\,\overline{\psi}(0,\omega^-,\mathbf{0}_\mathrm{T})
W(\omega^-,0)\gamma^+\frac{\lambda^a}{2}\psi(0) \right|
P\right\rangle_\mathrm{C},
\end{align}
where $T$ is the
time-ordering operator, $\psi$ is a quark
field, $\lambda^a$ is an SU(2)-flavor Pauli matrix, and the subscript C indicates that the vacuum
expectation value has been subtracted (in other words, only connected
contributions are included). The operator $W(\omega^-,0)$ is the Wilson line,
\begin{equation}
W(\omega^-,0) =
{\cal P}\exp\left[-ig_0\int_0^{\omega^-}\mathrm{d}y^-A^+_c(0,y^-,
\mathbf{0}_{\mathrm{T}})T_c\right],
\end{equation}
with ${\cal P}$ the path-ordering operator, $g_0$ the QCD bare coupling,
and $A^\alpha = A^\alpha_c T_c$ the $SU(3)$ gauge potential with generator
$T_c$ (summation over color index $c$ is implicit). The target
state, $|P\rangle$, is a spin-averaged, exact momentum eigenstate with
relativistic normalization
\begin{equation}
\langle P' | P \rangle =
(2\pi)^32P^+\delta\left(P^+-P^{\prime\,+}\right)\delta^{(2)}
\left(\mathbf{P}_\mathrm{T} - \mathbf{P}'_\mathrm{T}\right).
\end{equation}

\subsection{Smeared quasidistributions}

One of the central challenges for nonperturbative calculations of quasidistributions is a power divergence associated with the length of the Wilson line in the presence of the lattice regulator. A number of approaches have been proposed to remove this power divergence \cite{Chen:2016fxx,Ishikawa:2016znu,Chen:2017mzz,Alexandrou:2017huk,Orginos:2017kos}, including circumventing the issue completely by using alternative operators \cite{Aglietti:1998ur,Abada:2001if,Detmold:2005gg,Braun:2007wv,Bali:2017gfr}. In Ref.~\cite{Monahan:2016bvm}, we constructed a finite matrix element to extract the continuum quasidistribution from lattice calculations by smearing both the fermion and gauge fields via the gradient flow
\cite{Narayanan:2006rf,Luscher:2011bx,Luscher:2013cpa}. For a more complete discussion of the gradient flow, I refer the reader to the recent reviews
\cite{Luscher:2013vga,Ramos:2015dla}.

At finite flow time, the gradient flow guarantees that the matrix elements constructed from flowed fields is UV finite. Provided one holds the flow time fixed in physical units, the lattice regulator can be removed and one can extract a continuum quasidistribution that is a function of the flow time. This smeared quasidistribution can then be related to the light-front distributions through a perturbative matching relation \cite{Monahan:2016bvm}, analogous to the relation that holds for the ordinary (unsmeared) quasidistribution. Here, I study, at one loop in perturbation theory, the matrix element that defines the smeared quasidistribution. This result is necessary to relate the smeared matrix element to the corresponding matrix element in the $\ms$ scheme.

The continuum smeared matrix element is finite but, because of the presence of the flow-time scale, is expected to contain a power divergence as the flow time is taken to zero \cite{Monahan:2016bvm}. This contribution must be removed nonperturbatively. The gradient flow provides a continuous parameter, in contrast to the lattice regulator, which should provide a better lever through which to study the power-divergent contribution. At one loop in perturbation theory, this divergence manifests as a term linear in the length of the Wilson line, an expectation confirmed by the results presented in Sec.~\ref{sec:pertth}.

I denote the ringed fermion fields \cite{Makino:2014taa,Hieda:2016lly} at flow time $\tau$ by $\overline{\chi}(x;\tau)$ and $\chi(x;\tau)$, and the corresponding Wilson line
at the same flow time, constructed from the smeared gauge fields $B_\alpha(x;\tau)$, by ${\bf\cal W}(x_1,x_2;\tau)$. I start with the connected matrix element
\begin{align}\label{eq:hdef}
{} & \hm^{(s)}\left(\frac{n^2}{\tau},n\cdot P,
\sqrt{\tau}\Lambda_{\mathrm{QCD}},\sqrt{\tau}M_\mathrm{N}\right) = \nonumber\\
{} & \quad
\frac{1}{2P}\left\langle
P \left| \overline{\chi}(n;\tau)
{\bf\cal W}(n,0;\tau)\gamma_\alpha\frac{\lambda^a}{2}\chi(0;\tau)\right|
P\right\rangle_\mathrm{C},
\end{align}
which, being dimensionless, depends only on dimensionless combinations of scales. Here, $n$ is a 4-vector that is usually taken to be in the $z$-direction, $n = (\mathbf{0},z,0)$; the dependence of the matrix element on 4-vectors is constrained by Euclidean SO(4) invariance \cite{Radyushkin:2016hsy}. I note that the flow time has units of length squared. The ringed fermion fields require no wave function renormalization and this smeared matrix element is finite provided the flow time, $\tau$, is nonzero and fixed in physical units,
because correlation functions constructed from smeared fields are finite \cite{Luscher:2011bx,Luscher:2013cpa}. Divergences will appear in the
limit of vanishing flow time and the matrix element will then require renormalization. The Lorentz index $\alpha$ can take any value from 1 to 4, but the choice $\alpha = 4$ removes some higher-twist contamination \cite{Radyushkin:2016hsy} and avoids mixing, at least for the unpolarized flavor-nonsinglet distribution, at finite lattice spacing \cite{Constantinou:2017sej}. In lattice calculations, it is common to choose the spatial momentum of the nucleon state to be $\mathbf{P} = (0,0,P_z)$, in which case the index $\alpha$ is restricted to be $\alpha = \{3,4\}$.

The smeared quasidistribution \cite{Monahan:2016bvm} is then defined, assuming $n = (\mathbf{0},z,0)$, as
\begin{align}\label{eq:qpdfdef}
q^{\,(s)}{} & \left(\xi,\sqrt{\tau}P_z,\sqrt{\tau}\Lambda_{\mathrm{QCD}},
\sqrt{\tau}M_\mathrm{N}\right)
=
\int_{-\infty}^\infty \frac{\mathrm{d}z}{2\pi} e^{i\xi z
P_z} P_z\nonumber \\
{} & \qquad \times h^{(s)}\left(\frac{n^2}{\tau},n\cdot P,\sqrt{\tau}\Lambda_{\mathrm{QCD}},
\sqrt{\tau}M_\mathrm{N}\right),
\end{align}
where $\xi$ is a dimensionless parameter that in Euclidean space should be viewed as a (dimensionless) Fourier-conjugate variable.

\subsection{Relating smeared quasidistributions and PDFs}

The smeared quasidistribution can be directly related to the light-front distribution by \cite{Monahan:2016bvm}
\begin{align}
{} & q^{\,(s)}\left(x,\sqrt{\tau} \Lambda_{\mathrm{QCD}}, \sqrt{\tau}P_z\right) =  \nonumber \\
{} & \; \int_{-1}^{1}
\frac{d\xi}{\xi}\, \widetilde{Z}\left(\frac{x}{\xi},\sqrt{\tau}\mu,
\sqrt{\tau}P_z\right)
f(\xi,\mu) +
{\cal O}(\sqrt{\tau}\Lambda_{\mathrm{QCD}}),
\end{align}
provided
\begin{equation}\label{eq:scales_tau}
\Lambda_{QCD},M_N\ll P_z \ll \tau^{-1/2}.
\end{equation}
The quasidistribution and the light-front PDF have the same IR structure \cite{Ma:2014jla,Ma:2014jga,Carlson:2017gpk,Briceno:2017cpo}, so that the matching kernel can be determined in perturbation theory. In practice, it is simplest to relate the smeared quasidistribution to the quasidistribution at zero flow time first and then match this (unsmeared) quasidistribution to the light-front PDF \cite{Xiong:2013bka,Ji:2015jwa}.

I show the diagrams representing the tree-level and one-loop contributions to the smeared quasidistribution in generalized Feynman gauge (with $\alpha = \lambda = 1$, where $\alpha$ is the usual gauge-fixing parameter and $\lambda$ is a parameter introduced to fix the gauge of the smeared gauge fields \cite{Luscher:2010iy,Luscher:2011bx}) in Figs.~\ref{fig:sqpdf_tree} and \ref{fig:sqpdf1}, respectively. These diagrams correspond to the ``quark-in-quark'' distributions at one loop in perturbation theory. Although the quasidistribution can be more easily evaluated in axial gauge, this gauge cannot be consistently generalized to arbitrary flow time.
\begin{figure}
\centering
\includegraphics[width =.25\textwidth,keepaspectratio=true]{images/tree.pdf}
\caption{Tree-level contribution to the smeared nonsinglet quasidistribution. The gray circles indicate fermion fields at finite flow time $\tau$ and solid lines represent fermion propagators.}
\label{fig:sqpdf_tree}
\end{figure}
\begin{figure}
\centering
\includegraphics[width =0.4\textwidth]{images/qmatch1.pdf}\\
(a)\hspace*{0.12\textwidth}(b)\hspace*{0.12\textwidth}(c)\\[8pt]
\includegraphics[width =0.42\textwidth]{images/qmatch2.pdf}\\
(d)\hspace*{0.12\textwidth}(e)\hspace*{0.12\textwidth}(f)\\[8pt]
\includegraphics[width =0.48\textwidth]{images/qmatch_zpsi.pdf}\\
(g)\hspace*{0.1\textwidth}(h)\hspace*{0.1\textwidth}(i)\hspace*{0.1\textwidth}(j)
\caption{ Diagrams representing the one-loop contributions to the smeared nonsinglet quasidistribution. The gray circles indicate fermion fields at finite flow time $\tau$; solid lines represent fermion propagators; double lines are fermion flow kernels; and open squares are flow vertices at arbitrary flow time $\tau_1$. The diagrams labeled (a), (d), (f), and (g) are in one-to-one correspondence with the diagrams that contribute to the unsmeared quasidistribution in the $\ms$ scheme, for which the fermion fields that appear in the extended operator are replaced by fermion fields at vanishing flow time.}
\label{fig:sqpdf1}
\end{figure}
In principle, it is possible to evaluate these nine diagrams and determine the matching kernel directly at one loop in perturbation theory. In practice, however, the corresponding integrals cannot be evaluated analytically, and it is simpler to use a procedure similar to that employed in Refs.~\cite{Ishikawa:2016znu,Constantinou:2017sej} and match at a fixed value of the Wilson-line length, $z$. This procedure, which corresponds to matching at the level of the matrix element itself, rather than dealing with the quasidistributions, is preferable largely because the matrix elements do not have any dependence on $\xi$. In the rest of this work, I focus on the matrix element, $\hm^{(s)}$, itself.
